% Part opener. Structural figures get their own number prefix so they
% do not inherit a misleading chapter number.
\structuralfigures{P2.}

\vspace*{0.4cm}
This part and the next are half the exam between them. Domain 2, Switching and
Network Access, is 25\% on its own, and it is the domain where the blueprint
stops asking you to explain things and starts asking you to \emph{configure}
them: three of its five topics begin with that word.

It is also the part where a HEC-trained student is furthest from the exam. No
published HEC networking course outline in any edition names a VLAN, a trunk, an
EtherChannel or a spanning tree. Everything in these seven chapters is new
ground, and all of it is typed rather than read.

\ccnafigh{%
\begin{tikzpicture}[font=\scriptsize]
  \tikzset{
    stg/.style={rounded corners=3pt, draw=#1, fill=#1!8, text=#1!45!black,
                line width=1pt, text width=3.15cm, align=center,
                minimum height=1.55cm, font=\scriptsize\bfseries},
    fl/.style={-{Stealth[length=2.4mm]}, line width=0.9pt, neutral!70}}

  \node[stg=dom2] (a) at (0,0)     {\faIcon{th-large}\\[2pt]8--9. VLANs\\
      \footnotesize\mdseries divide the switch};
  \node[stg=dom2] (b) at (4.0,0)   {\faIcon{exchange-alt}\\[2pt]10. Trunks\\
      \footnotesize\mdseries carry them between switches};
  \node[stg=dom2] (c) at (8.0,0)   {\faIcon{link}\\[2pt]11. EtherChannel\\
      \footnotesize\mdseries bundle the links};
  \node[stg=dom2] (d) at (12.0,0)  {\faIcon{sitemap}\\[2pt]12. Rapid PVST+\\
      \footnotesize\mdseries survive the loops};
  \node[stg=dom1] (e) at (4.0,-2.7) {\faIcon{clipboard-list}\\[2pt]13. CDP, LLDP\\
      \footnotesize\mdseries prove the diagram};
  \node[stg=accent] (f) at (8.0,-2.7) {\faIcon{bug}\\[2pt]14. Troubleshooting\\
      \footnotesize\mdseries when it does not work};

  \draw[fl] (a) -- (b);
  \draw[fl] (b) -- (c);
  \draw[fl] (c) -- (d);
  \draw[fl] (d.south) -- ++(0,-0.6) -| (e.north);
  \draw[fl] (e) -- (f);
\end{tikzpicture}}%
{Part~II. Each chapter adds one capability to the same switched network, and the
last two are about proving it works and finding out why it does not.}{p2map}

\begin{examnote}[Where the marks are in domain 2]
Topic \tp{2.1} (switch-to-switch and switch-to-router connectivity) and
\tp{2.2} (edge-host port attributes) are both \emph{Configure} topics and between
them cover trunks, EtherChannel, SVIs, VLANs and PoE. Topic \tp{2.5} is
\emph{Configure} as well --- Rapid PVST+ with PortFast and the guards. Topic
\tp{2.3} asks you to \emph{validate} documentation with CDP and LLDP, and
\tp{2.4} to \emph{troubleshoot} with \cmd{show} commands, ping, traceroute and
capture output.

Four of the five are hands-on. Read this part with a simulator open.
\end{examnote}
